\clearpage
\FrontTitle{ABSTRAK}
\begin{center}
{\bfseries\large \JudulLaporan\par
\bfseries\large \TempatKP\par}
\vspace{0.1cm}
{Oleh\par
\bfseries\large Naufarrel Zhafif Abhista\par
\bfseries\large NIM: \NIM\par}
\end{center}
\vspace{1.1cm}

PT Radya Anugrah Digital (Radya Digital) adalah perusahaan pengembang perangkat lunak di Bandung yang sedang mengerjakan aplikasi Komatsu ESS untuk klien, yaitu Komatsu. Aplikasi itu perlu diuji secara terencana supaya berfungsi sesuai \textit{requirement} klien dan bebas dari \textit{bug} berarti sebelum diserahkan. KP ini dilaksanakan pada 22 Juni sampai 22 September 2026 dengan tujuan memastikan kualitas aplikasi Komatsu ESS lewat kegiatan \textit{Quality Assurance} (QA). Pekerjaan dilakukan bersama seorang rekan sesama mahasiswa KP dengan tahapan penyusunan draf \textit{test script}, \textit{module testing}, \textit{integration testing}, dan \textit{regression testing}. \textit{Bug} yang ditemukan di setiap tahap dilaporkan ke \textit{developer} untuk diperbaiki, lalu perbaikannya diverifikasi ulang. Hasil KP berupa \textit{test script}, \textit{bug report}, dan laporan hasil pengujian yang dipakai \textit{developer} untuk meningkatkan kualitas aplikasi. \TEKNIS{tambahkan hasil kuantitatif, misalnya jumlah modul, \textit{test script}, dan \textit{bug}} Selain itu, KP ini melatih kerja sama tim, komunikasi dengan \textit{developer}, dan sikap profesional di lingkungan kerja.

Kata kunci: pengujian perangkat lunak, \textit{quality assurance}, \textit{regression testing}, \textit{integration testing}.
